% !TEX root = ../main.tex
%=====================================================================
%  扉页 / 内容提要 / 译者说明 / 目录
%  原文：G. W. Moore, SCGP 2017 讲义，第 1 页
%=====================================================================

\begin{titlepage}
\centering
\vspace*{2.2cm}

{\LARGE\bfseries 用物理方法研究四维流形的\\[0.35em]
唐纳森不变量与塞伯格--威滕不变量\par}

\vspace{1.0cm}

{\large 讲义\par}

\vspace{1.6cm}

{\Large Gregory W.\ Moore\par}
\vspace{0.5em}
{\normalsize NHETC and Department of Physics and Astronomy,\\
Rutgers University\\
\texttt{[邮箱隐去]}\par}

\vspace{2.2cm}

{\normalsize
Simons Center for Geometry and Physics 讲义系列\\
2017 年 3 月 22--24 日\\
（删节本曾于 2017 年 7 月 17--18 日在汉堡 Pre-String-Math 学校讲授）\par}

\vspace{1.2cm}
{\normalsize 中译本\par}

\vfill
{\small 原稿日期：2017 年 8 月 25 日\quad 原稿状态：仍在修订中（THESE NOTES ARE UNDER CONSTRUCTION）\par}
\vspace{0.6cm}
\end{titlepage}

%---------------------------------------------------------------------
\section*{内容提要}
\addcontentsline{toc}{section}{内容提要}

本讲义讨论四维流形拓扑中两个里程碑式的不变量理论——唐纳森（Donaldson）不变量与
塞伯格--威滕（Seiberg--Witten）不变量——之间的深刻联系，并系统展示如何用
$\mathcal{N}=2$ 超对称规范场论的量子场论方法来“推导”这一联系。

全篇的主线可以概括为三步。第一步（第 3 节）建立上同调拓扑场论的公理化框架，
即路径积分的 Mathai--Quillen 形式与局域化（localization）机制：一个具有幂零超对称
算子 $Q$ 的路径积分，在 $Q$-恰当形变下不变，因而精确局域化到 $Q$ 不动点方程的解集
（模空间）上，并化为该模空间上的上同调类交截数。第二步（第 4 节）指出，把
$\mathcal{N}=2$ 超对称杨--米尔斯理论做拓扑扭曲（取 $R$-对称丛等于自对偶形式丛，
$\omega_{R}=\omega^{+}$）之后，理论恰好落入上述框架，其 $Q$ 不动点方程就是反自对偶
瞬子方程 $F^{+}=0$，而可观测量的关联函数就是唐纳森多项式。第三步（第 5--9 节）是
全篇的技术核心：把紫外线（UV）理论映射到红外（IR）的低能有效理论（LEET）。
由于理论是拓扑的，可以把度规整体放大 $g\to tg$ 并取 $t\to\infty$，此时物理完全由
真空模空间（$u$ 平面）及其上的低能有效作用量决定；路径积分随之分解为
“库仑分支贡献（$u$ 平面积分）”与“希格斯分支贡献（塞伯格--威滕模空间上的积分）”
两部分之和：
\[
  Z^{\xi}_{\mathrm{DW}}(p,\Sigma)
  \;=\; Z^{\xi}_{u}(p,\Sigma)\;+\;Z^{\mathrm{IR},\mathrm{Higgs}}(p,\Sigma).
\]
对 $b^{+}_{2}(X)>1$ 且 $X$ 为塞伯格--威滕单型的流形，$u$ 平面积分消失，配分函数
完全由塞伯格--威滕不变量给出，这正是威滕 1994 年提出的著名猜想；而对 $b^{+}_{2}=1$
的情形，$u$ 平面积分与塞伯格--威滕不变量各自的穿墙跳变相互抵消，这一自洽性要求反过来
唯一地确定了低能有效作用量中原本未知的耦合函数 $C(u),P(u),E(u)$——这是全篇最漂亮的
论证。讲义末尾讨论该物理图像如何解释已知数学定理（后测），又如何给出新的数学预言
（超共形单型、广义诺特不等式、四维流形地理问题）。

%---------------------------------------------------------------------
\section*{译者说明}
\addcontentsline{toc}{section}{译者说明}

\begin{enumerate}[leftmargin=2.2em, itemsep=4pt]
\item \textbf{底本。}译本据作者 2017 年 8 月 25 日版 PDF（82 页，pdfTeX 生成）逐页译出，
      保留原文的章节结构、公式编号与全部脚注。原文为讲义手稿，作者自注“仍在修订中”，
      正文中保留了若干未完成的占位符（如 \texttt{[CITE: \dots]}、大写提醒句，
      以及以 $\clubsuit$ 括起的批注）。这些作者自注在译本中以脚注形式标出，冠以“〔原稿边注〕”，
      以区别于译者批注。

\item \textbf{公式编号。}全部公式沿用原文编号（如 (4.59)、(7.7)、(A.14)），
      不重新编号，以便与原文及文献对照。原文中被作者略去未编号之处，
      译本亦不擅自补号。

\item \textbf{译者批注。}凡涉及符号约定陷阱、作者略去的证明、隐含假设、
      以及与前后文的呼应，译本在相应段落的 LaTeX 源码中以
      \verb|% [Tutor Note: ...]| 的形式插入批注。批注只出现在源码中，
      不进入排版结果，以便读者在编辑 \texttt{.tex} 文件时直接看到。
      全书共若干处，均以“【易错】”“【省略】”“【约定】”等起首标明性质。

\item \textbf{本节重点。}每节开头附“本节重点（Key Takeaways）”栏目，
      提炼该节的物理图像、所用数学工具与核心结论。

\item \textbf{图表。}原 PDF 中\textbf{并无任何实际图形}：作者仅在正文中留下
      十余处图位说明（如 \texttt{FIGURE: CHAMBERS}），以及一处带图号的图题
      （\texttt{FIGURE1: SOME KNOWN BOUNDS}，第 70 页），而图形本身从未插入。
      为忠实于原文，译本不伪造任何图形，而是以“图位占位框”重现原文的图位，
      并将图题译为中文，同时注明其在原 PDF 中的页码。

\item \textbf{译名。}人名、术语依中文学界通行译法；首次出现时括注原文。
      主要译名对照见下表。凡有歧义者（如 \emph{chamber} 译“腔室”、
      \emph{dyon} 译“双荷子”），均以括注保留英文原词。
\end{enumerate}

\vspace{0.6em}
\noindent\textbf{主要术语对照}\quad
规范场论 gauge theory；瞬子 instanton；反自对偶 anti-self-dual (ASD)；
模空间 moduli space；相交形式 intersection form；反常 anomaly；鬼数 ghost number；
拓扑扭曲 topological twisting；等变上同调 equivariant cohomology；局域化 localization；
低能有效理论 LEET；库仑分支 Coulomb branch；希格斯分支 Higgs branch；预势 prepotential；
对偶标架 duality frame；单值性 monodromy；穿墙 wall-crossing；单型 simple type；
基本类 basic class；爆破 blow-up；地理问题 geography problem。

\vspace{1.2em}
\noindent\textbf{编译方式}\quad 本源码以 Xe\LaTeX{} 编译（两遍以生成目录与交叉引用）：
\begin{verbatim}
    xelatex main.tex
    xelatex main.tex
\end{verbatim}
中文支持由 \texttt{ctex} 宏包提供（默认 \texttt{fontset=fandol}，
以保证跨平台可编译；在 Windows 上可改为 \texttt{fontset=windows} 以使用宋体/黑体）。

\clearpage
\tableofcontents
\clearpage
